\chapter{Geometric and height infrastructure}
\label{chap:basic}

% archon:covers MR4276287UniformityInMordellLangForCurves/Basic.lean

This chapter packages the geometric and height-theoretic input used by the
later counting arguments.

\section{Introduction}

\begin{lemma}[Height comparison on the moduli side]
  \label{lem:moduli_height_comparison}
  \lean{MR4276287UniformityInMordellLangForCurves.Basic.moduliHeightComparison}
  \uses{def:ambient_setup,def:torelli_map,lem:torelli_quasi_finite}
  % SOURCE: Introduction of \cite{MR4276287-mordell-lang-curves} (read from references/MR4276287-mordell-lang-curves-source/intro.tex)
  % SOURCE QUOTE: "If $k=\IQbar$ and as described in $\mathsection\ref{SubsectionHeight}$ we obtain two height functions, the restriction of the Weil height $h\colon \overline{S}(\IQbar)\rightarrow \IR$ and the N\'eron--Tate height $\hat h_{\cA} \colon \cA(\IQbar)\rightarrow \IR$."
  % SOURCE QUOTE PROOF: The paper compares these heights after fixing a projective embedding on the moduli side and reading the induced height along the Torelli map.
  \textit{Source: Introduction of \cite{MR4276287-mordell-lang-curves} (read from references/MR4276287-mordell-lang-curves-source/intro.tex).}
  The chosen projective height on the moduli space and the height induced by
  \(\mathbb{A}_g\) are comparable up to affine linear error terms along the
  Torelli map.
\end{lemma}

% SOURCE QUOTE PROOF: "If $k=\IQbar$ and as described in $\mathsection\ref{SubsectionHeight}$ we obtain two height functions, the restriction of the Weil height $h\colon \overline{S}(\IQbar)\rightarrow \IR$ and the N\'eron--Tate height $\hat h_{\cA} \colon \cA(\IQbar)\rightarrow \IR$." (read from references/MR4276287-mordell-lang-curves-source/intro.tex)
\begin{proof}
  \uses{def:ambient_setup, thm:silverman_tate}
  Fix a projective embedding \(\overline{\mathbb{A}_g} \subset \mathbb{P}^m\) inducing
  the Weil height \(h\) on \(\overline{\mathbb{A}_g}(\overline{\mathbb{Q}})\).
  The Torelli morphism \(\tau \colon \mathbb{M}_g \to \mathbb{A}_g\) is quasi-finite,
  so by the Height Machine (functoriality of height under morphisms of projective varieties,
  \cite[Lemma~4]{Silverman:heightest11}) there exist constants \(c' > 0\) and \(c'' \ge 0\)
  such that \(h_{\overline{\mathbb{A}_g},\mathcal{M}}(\tau(s)) \le c' h(s) + c''\)
  and \(h(s) \le c' h_{\overline{\mathbb{A}_g},\mathcal{M}}(\tau(s)) + c''\)
  for all \(s \in \mathbb{M}_g(\overline{\mathbb{Q}})\).
  On each fiber \(\mathcal{A}_s\) the Néron--Tate height \(\hat{h}_{\mathcal{A}}\) and
  the naive projective height \(h\) differ by a term bounded linearly in
  \(h_{\overline{S}}(\pi(P))\) by \cref{thm:silverman_tate}.
  Composing these two comparisons yields the stated affine linear comparability.
\end{proof}

\section{Betti map and Betti form}

\subsection{Betti map for the universal abelian variety}

\begin{proposition}[Betti map for the universal abelian variety]
  \label{prop:betti_map}
  \lean{MR4276287UniformityInMordellLangForCurves.Basic.bettiMap}
  \uses{def:ambient_setup}
  % SOURCE: Section 2 of \cite{MR4276287-mordell-lang-curves} (read from references/MR4276287-mordell-lang-curves-source/BettiMapForm.tex)
  % SOURCE QUOTE: "There exist an open neighborhood \(\Delta\) of \(s_0\) in \(S^{\mathrm{an}}\), and a map \(b_{\Delta}\colon \cA_{\Delta} := \pi^{-1}(\Delta) \rightarrow \mathbb T^{2g}\), called the Betti map, with the following properties."
  % SOURCE QUOTE PROOF: This is the local Betti-map construction on the universal abelian family.
  \textit{Source: Section 2 of \cite{MR4276287-mordell-lang-curves} (read from references/MR4276287-mordell-lang-curves-source/BettiMapForm.tex).}
  Locally on the base, the analytic family of abelian varieties admits a Betti
  map to \(\mathbb{T}^{2g}\) whose restriction to each fiber is a group
  isomorphism, whose fibers are complex-analytic, and whose product with the
  projection to the base is a real-analytic isomorphism.
\end{proposition}

% SOURCE QUOTE PROOF: "Our proof of Proposition~\ref{PropBettiMap} follows the construction in \cite[$\mathsection$3-$\mathsection$4]{GaoBettiRank}. We divide it into several steps. We start to prove Proposition~\ref{PropBettiMap} for $S = \mathbb A_g$, the moduli space of principally polarized abelian variety of dimension $g$ with level-$\ell$-structure; it is a fine moduli space. Let $\pi^{\mathrm{univ}} \colon \mathfrak{A}_g \rightarrow \mathbb A_g$ be the universal abelian variety. The universal covering $\mathfrak{H}_g \rightarrow \an{\mathbb A}_g$, where $\mathfrak{H}_g$ is the Siegel upper half space, gives a polarized family of abelian varieties $\cA_{\mathfrak{H}_g} \rightarrow \mathfrak{H}_g$." (read from references/MR4276287-mordell-lang-curves-source/BettiMapForm.tex)
\begin{proof}
  \uses{def:ambient_setup}
  First treat the case \(S = \mathbb{A}_g\). The universal covering
  \(\mathfrak{H}_g \to \mathbb{A}_g^{\mathrm{an}}\) by the Siegel upper half-space
  gives a family \(\mathcal{A}_{\mathfrak{H}_g} \to \mathfrak{H}_g\).
  For the universal covering \(u \colon \mathbb{C}^g \times \mathfrak{H}_g \to \mathcal{A}_{\mathfrak{H}_g}\)
  and each \(Z \in \mathfrak{H}_g\), the kernel of \(u|_{\mathbb{C}^g \times \{Z\}}\)
  is \(\mathbb{Z}^g + Z\mathbb{Z}^g\). The map
  \(\mathbb{C}^g \times \mathfrak{H}_g \to \mathbb{R}^{2g}/\mathbb{Z}^{2g}\)
  defined by the inverse of \((a,b,Z) \mapsto (a + Zb, Z)\) followed by projection
  descends to a real-analytic map \(b^{\mathrm{univ}} \colon \mathcal{A}_{\mathfrak{H}_g} \to \mathbb{T}^{2g}\).
  For any \(s_0 \in \mathbb{A}_g(\mathbb{C})\), a contractible open neighborhood
  \(\Delta\) of \(s_0\) in \(\mathbb{A}_g^{\mathrm{an}}\) identifies
  \(\mathfrak{A}_{g,\Delta}\) with \(\mathcal{A}_{\mathfrak{H}_g,\Delta'}\) for some
  \(\Delta' \subset \mathfrak{H}_g\); the composite \(b_\Delta \colon \mathfrak{A}_{g,\Delta}
  \cong \mathcal{A}_{\mathfrak{H}_g,\Delta'} \to \mathbb{T}^{2g}\) satisfies all three
  properties.
  The general case (for \(\mathcal{A} \to S\) satisfying {\tt (Hyp)}) follows from
  the fine moduli property of \(\mathbb{A}_g\): the Cartesian diagram
  \(\mathcal{A} \to \mathfrak{A}_g\) over \(S \xrightarrow{\iota_S} \mathbb{A}_g\)
  gives \(b_\Delta = b_{\Delta_0} \circ \iota\), which inherits all required properties.
\end{proof}

\subsection{Betti form for the universal abelian variety}

\begin{proposition}[Betti form]
  \label{prop:betti_form}
  \lean{MR4276287UniformityInMordellLangForCurves.Basic.bettiForm}
  \uses{prop:betti_map}
  % SOURCE: Section 2 of \cite{MR4276287-mordell-lang-curves} (read from references/MR4276287-mordell-lang-curves-source/BettiMapForm.tex)
  % SOURCE QUOTE: "There exists a closed \((1,1)\)-form \(\omega\) on \(\an{\cA}\), called the Betti form, such that the following properties hold."
  % SOURCE QUOTE PROOF: This is the companion form attached to the same universal family, and it controls Betti rank through positivity.
  \textit{Source: Section 2 of \cite{MR4276287-mordell-lang-curves} (read from references/MR4276287-mordell-lang-curves-source/BettiMapForm.tex).}
  There is a closed semi-positive \((1,1)\)-form on the analytification of the
  universal family which scales by \(N^2\) under multiplication-by-\(N\), and
  whose top wedge detects the maximal real rank of the Betti map on smooth
  loci.
\end{proposition}

% SOURCE QUOTE PROOF: "Define $\hat{\omega}^{\mathrm{univ}}:=\sqrt{-1}\partial \overline{\partial} \left( 2 (\mathrm{Im}w)^{\!^{\intercal}} (\mathrm{Im} Z)^{-1} (\mathrm{Im}w) \right)$. Then $\hat{\omega}^{\mathrm{univ}}$ is a closed semi-positive $(1,1)$-form on $\IC^g\times\mathfrak{H}_g$." (read from references/MR4276287-mordell-lang-curves-source/BettiMapForm.tex)
\begin{proof}
  \uses{prop:betti_map}
  On the universal cover \(\mathbb{C}^g \times \mathfrak{H}_g\), define
  \(\hat{\omega}^{\mathrm{univ}} := \sqrt{-1}\,\partial\bar{\partial}
  \bigl(2(\mathrm{Im}\,w)^\intercal (\mathrm{Im}\,Z)^{-1}(\mathrm{Im}\,w)\bigr)\).
  A direct computation shows
  \(\hat{\omega}^{\mathrm{univ}} = 2\,d^\intercal a \wedge db\)
  in the real coordinates \((a,b,Z)\) defined by \(w = a + Zb\),
  from which semi-positivity follows because \(\mathrm{Im}(Z)\) is positive definite.
  This form is invariant under the action of \(P_{2g,\mathrm{a}}(\mathbb{R}) = V_{2g}(\mathbb{R}) \rtimes \mathrm{Sp}_{2g}(\mathbb{R})\)
  on \(\mathcal{X}_{2g,\mathrm{a}} = \mathbb{R}^{2g} \times \mathfrak{H}_g\),
  hence it descends to a semi-positive \((1,1)\)-form \(\omega^{\mathrm{univ}}\) on
  \(\mathfrak{A}_g^{\mathrm{an}}\).
  The scaling \([N]^*\omega^{\mathrm{univ}} = N^2\omega^{\mathrm{univ}}\) holds because
  the lift \(\widetilde{N} \colon (w,Z) \mapsto (Nw, Z)\) scales \(\hat{\omega}^{\mathrm{univ}}\)
  by \(N^2\).
  The Betti form on \(\mathcal{A}^{\mathrm{an}}\) is \(\omega := \iota^*\omega^{\mathrm{univ}}\),
  inheriting semi-positivity and the scaling property.
  The equivalence between \(\omega|_{X^{\mathrm{sm}}}^{\wedge d} \not\equiv 0\)
  and maximal Betti rank \(2d\) on smooth points is proved by showing that the
  Hermitian form of \(\omega|_{X^{\mathrm{sm}}}\) is degenerate precisely when
  every point of \(X^{\mathrm{sm}} \cap \mathcal{A}_\Delta\) is a non-isolated point of
  a level set of \(b_\Delta\).
\end{proof}

\section{Intersection theory and height inequality on the total space}

\subsection{The graph construction}

\begin{lemma}[Graph construction for the height inequality]
  \label{lem:graph_construction}
  \lean{MR4276287UniformityInMordellLangForCurves.Basic.graphConstruction}
  \uses{def:ambient_setup}
  % SOURCE: Section 4 of \cite{MR4276287-mordell-lang-curves} (read from references/MR4276287-mordell-lang-curves-source/HtIneq.tex)
  % SOURCE QUOTE: "For each \(N \in \IN\), we work with \(X_N \subseteq \cA\times_S \cA\), the graph of \([N] \colon \cA \rightarrow \cA\), and its Zariski closure \(\overline{X_N}\) in \(\overline{\cA}\times_{\overline{S}}\overline{\cA}\)."
  % SOURCE QUOTE PROOF: This is the projective closure on which the auxiliary inequality is evaluated.
  \textit{Source: Section 4 of \cite{MR4276287-mordell-lang-curves} (read from references/MR4276287-mordell-lang-curves-source/HtIneq.tex).}
  For each \(N\), the graph of multiplication-by-\(N\) on the family produces a
  projective closure \(\overline{X_N}\) on which the relevant heights become the
  heights associated to the line bundles \(\mathcal F\) and \(\mathcal M\).
\end{lemma}

% SOURCE QUOTE PROOF: "For each $N \in \IN$, we work with $X_N \subseteq \cA\times_S \cA$, the graph of $[N] \colon \cA \rightarrow \cA$, and its Zariski closure $\overline{X_N}$ in $\overline{\cA}\times_{\overline{S}}\overline{\cA} \subseteq \IP_k^n \times \IP_k^n \times \IP_k^m$. The point $P$ gives rise to a point $P_N \in \overline{X_N}$. Consider the line bundles $\cF = \cO(0,1,1)|_{\overline{X_N}}$ and $\cM = \cO(0,0,1)|_{\overline{X_N}}$." (read from references/MR4276287-mordell-lang-curves-source/HtIneq.tex)
\begin{proof}
  \uses{def:ambient_setup}
  For each \(N \in \mathbb{N}\), define \(X_N \subset \mathcal{A} \times_S \mathcal{A}\)
  to be the graph of the multiplication morphism \([N] \colon \mathcal{A} \to \mathcal{A}\).
  Since \(\mathcal{A} \subset \mathbb{P}^n \times S\) and \(S \subset \mathbb{P}^m\),
  we have \(\mathcal{A} \times_S \mathcal{A} \subset \mathbb{P}^n \times \mathbb{P}^n \times S\).
  Let \(\overline{X_N}\) be the Zariski closure of \(X_N\) inside
  \(\overline{\mathcal{A}} \times_{\overline{S}} \overline{\mathcal{A}} \subset
  \mathbb{P}^n \times \mathbb{P}^n \times \mathbb{P}^m\).
  Define the line bundles
  \(\mathcal{F} = \mathcal{O}(0,1,1)|_{\overline{X_N}}\) and
  \(\mathcal{M} = \mathcal{O}(0,0,1)|_{\overline{X_N}}\).
  A point \(P \in X(\overline{\mathbb{Q}})\) gives rise to a point
  \(P_N = (P, [N]P) \in X_N(\overline{\mathbb{Q}})\), and the height functions
  \(h_{\overline{X_N},\mathcal{F}}\) and \(h_{\overline{X_N},\mathcal{M}}\)
  attached by the Height Machine to \(\mathcal{F}\) and \(\mathcal{M}\)
  represent \(h([N]P)\) and \(h(\pi(P))\) respectively, up to bounded error.
\end{proof}

\subsection{Proof of Proposition~\ref{PropAuxHtIneq}}

\begin{proposition}[Auxiliary height inequality]
  \label{prop:aux_ht_ineq}
  \lean{MR4276287UniformityInMordellLangForCurves.Basic.auxHeightInequality}
  \uses{prop:betti_form,def:nondegenerate_app,lem:graph_construction,thm:siuCriterion,lem:nefFromBetti,lem:siu_intersection_lower,lem:siu_intersection_upper,thm:arith_bezout}
  % SOURCE: Section 4 of \cite{MR4276287-mordell-lang-curves} (read from references/MR4276287-mordell-lang-curves-source/HtIneq.tex)
  % SOURCE QUOTE: "Suppose \(\an{X}\) contains a smooth point at which \(\omega|_{\sman{X}}^{\wedge d} >0\). Then there exists a constant \(c_1 > 0\) ... \(h([N]P) \ge c_1 N^2 h(\pi(P)) - c_2(N)\)."
  % SOURCE QUOTE PROOF: The proof studies the graph of multiplication-by-\(N\), uses the Betti form to force positivity there, and then reads off a height lower bound.
  \textit{Source: Section 4 of \cite{MR4276287-mordell-lang-curves} (read from references/MR4276287-mordell-lang-curves-source/HtIneq.tex).}
  The graph construction turns the estimate into a comparison of intersection
  numbers on the projective closure \(\overline{X_N}\). Siu's criterion then
  converts the positivity of the Betti form into a lower bound for the
  projective height of \([N]P\) in terms of the height of the base point.
\end{proposition}

% SOURCE QUOTE PROOF: "Both $\cF$ and $\cM$ are nef line bundles. Thus a criterion of bigness by Siu \cite[Theorem~2.2.15]{PosAlgGeom}, states that $\cF^{\otimes q} \otimes \cM^{\otimes -p N^2}$ is big if $(\cF^{\cdot d}) > d c_1 (\cM^{\otimes N^2} \cdot \cF^{\cdot (d-1)})$. Note that $(\cM^{\otimes N^2} \cdot \cF^{\cdot (d-1)}) = N^2 (\cM \cdot \cF^{\cdot (d-1)})$ by multi-linearity of intersection numbers. Thus our task becomes comparing two intersection numbers." (read from references/MR4276287-mordell-lang-curves-source/HtIneq.tex)
\begin{proof}
  \uses{prop:betti_form, lem:graph_construction, def:nondegenerate_app}
  By hypothesis there exists a smooth point of \(X^{\mathrm{an}}\) where
  \(\omega|_{X^{\mathrm{sm,an}}}^{\wedge d} > 0\). Using the graph construction
  from \cref{lem:graph_construction}, working on the projective closure
  \(\overline{X_N} \subset \mathbb{P}^n \times \mathbb{P}^n \times \mathbb{P}^m\)
  with line bundles \(\mathcal{F} = \mathcal{O}(0,1,1)|_{\overline{X_N}}\) and
  \(\mathcal{M} = \mathcal{O}(0,0,1)|_{\overline{X_N}}\),
  the height inequality is equivalent to showing \(\mathcal{F}^{\otimes q} \otimes \mathcal{M}^{\otimes -qc_1 N^2}\)
  is big for suitable \(q, c_1 > 0\) independent of \(N\).
  The positivity of \(\omega|_{X^{\mathrm{sm}}}^{\wedge d}\) at a smooth point,
  together with the scaling \([N]^*\omega = N^2\omega\) from \cref{prop:betti_form},
  gives a lower bound \((\mathcal{F}^{\cdot d}) \ge \kappa N^{2d}\) for some \(\kappa > 0\)
  independent of \(N\), via a Lelong-type integration argument on \(X^{\mathrm{sm,an}}\).
  An upper bound \((\mathcal{M} \cdot \mathcal{F}^{\cdot(d-1)}) \le c N^{2(d-1)}\)
  follows by expressing \([N] = [2^l]\) through bihomogeneous polynomials of
  controlled bidegree and applying a Fulton-type positivity argument.
  Choosing \(c_1\) satisfying \(0 < c_1 \cdot c \cdot d < \kappa\), Siu's
  bigness criterion \cite[Theorem~2.2.15]{PosAlgGeom} gives the desired big
  line bundle, from which the height inequality follows by the Height Machine.
\end{proof}

\section{Proof of the height inequality Theorem~\ref{ThmHtInequality}}

\begin{theorem}[Height inequality under the paper's hypotheses]
  \label{thm:ht_inequality}
  \lean{MR4276287UniformityInMordellLangForCurves.Basic.heightInequality}
  \uses{prop:aux_ht_ineq,thm:silverman_tate,lem:graph_construction}
  % SOURCE: Section 5 of \cite{MR4276287-mordell-lang-curves} (read from references/MR4276287-mordell-lang-curves-source/NTbase.tex)
  % SOURCE QUOTE: "There exists a constant \(c_1>0\) ... \(\hat{h}_{\mathcal{A}}(P) \ge c_1 N^2 h(\pi(P)) - c_2(N)\)."
  % SOURCE QUOTE PROOF: This is the canonical-height form of the auxiliary estimate after the Silverman--Tate comparison is inserted.
  \textit{Source: Section 5 of \cite{MR4276287-mordell-lang-curves} (read from references/MR4276287-mordell-lang-curves-source/NTbase.tex).}
  After one fixed power of \(2\) is chosen, the canonical height on the total
  space is bounded below by a positive multiple of the height on the base on a
  dense open subset of any non-degenerate subvariety.
\end{theorem}

% SOURCE QUOTE PROOF: "By the Theorem of Silverman-Tate, see \cite[Theorem~A]{Silverman} and Theorem~\ref{thm:silvermantate}, there exist a constant $c_0\ge 0$ such that $|\hat{h}_{\mathcal{A}}(P)-h(P)| \le c_0\max\{1,h(\pi(P))\}\le c_0\bigl(1+h(\pi(P))\bigr)$ for all $P\in \mathcal{A}(\IQbar)$. We use $\hat{h}_{\mathcal{A}}([N]P)=N^2 \hat{h}_{\mathcal{A}}(P)$, divide by $N^2$, and rearrange to get $\hat{h}_{\cA}(P) \ge \left(c_1 - \frac{c_0}{N^2} \right) h(\pi(P)) - \frac{c_2(N)+ c_0}{N^2}$." (read from references/MR4276287-mordell-lang-curves-source/NTbase.tex)
\begin{proof}
  \uses{prop:aux_ht_ineq, thm:silverman_tate, lem:graph_construction}
  The non-degeneracy hypothesis and \cref{prop:betti_form}(iii) give
  \(\omega|_{X^{\mathrm{sm,an}}}^{\wedge d} > 0\) at some smooth point,
  so \cref{prop:aux_ht_ineq} applies. It yields \(c_1 > 0\) and,
  for each power-of-two \(N\), a Zariski open dense \(U_N \subset X\)
  and constant \(c_2(N)\) with
  \(h([N]P) \ge c_1 N^2 h(\pi(P)) - c_2(N)\) on \(U_N(\overline{\mathbb{Q}})\).
  By \cref{thm:silverman_tate}, there is \(c_0 \ge 0\) with
  \(\bigl|\hat{h}_{\mathcal{A}}(P) - h(P)\bigr| \le c_0(1 + h(\pi(P)))\).
  Applying this to \([N]P\) and using \(\hat{h}_{\mathcal{A}}([N]P) = N^2\hat{h}_{\mathcal{A}}(P)\):
  \[\hat{h}_{\mathcal{A}}(P) \ge \Bigl(c_1 - \frac{c_0}{N^2}\Bigr) h(\pi(P)) - \frac{c_2(N) + c_0}{N^2}.\]
  Choose \(N\) the least power of \(2\) with \(N^2 \ge 2c_0/c_1\), so that
  \(c_1 - c_0/N^2 \ge c_1/2 > 0\). The Zariski open dense subset \(U_N\) and the
  resulting constants give the stated inequality.
\end{proof}

\section{Preparation for counting points}

\subsection{The universal family and non-degeneracy}

\begin{definition}[Ambient moduli and height data]
  \label{def:ambient_setup}
  \lean{MR4276287UniformityInMordellLangForCurves.Basic.ambientSetup}
  \uses{def:Mg_level,def:universal_curve,def:Ag_level,def:universal_av,def:torelli_map}
  % SOURCE: Section 6 of \cite{MR4276287-mordell-lang-curves} (read from references/MR4276287-mordell-lang-curves-source/SettingUp.tex)
  % SOURCE QUOTE: "There exists a universal curve $\mathfrak{C}_g$ over $\mathbb{M}_g$,  it is smooth and proper over $\mathbb{M}_g$ and its fibers are smooth curves of genus $g$."
  % SOURCE QUOTE PROOF: This section fixes the universal curve, the universal abelian variety, and the height package used throughout the chapter.
  \textit{Source: Section 6 of \cite{MR4276287-mordell-lang-curves} (read from references/MR4276287-mordell-lang-curves-source/SettingUp.tex).}
  Fix \(g \ge 2\). Let \(\mathbb{M}_g\) be the fine moduli space of smooth
  genus-\(g\) curves with level structure, let \(\mathfrak{C}_g \to \mathbb{M}_g\)
  be the universal curve, and let \(\mathfrak{A}_g \to \mathbb{A}_g\) be the
  universal principally polarized abelian variety with the chosen projective
  embeddings. Fix the ample line bundles and the associated height functions on
  the compactifications, together with the fiberwise Néron--Tate height on the
  universal abelian family.
\end{definition}

\begin{lemma}[Uniform degree bound]
  \label{lem:uniform_degree_bound}
  \lean{MR4276287UniformityInMordellLangForCurves.Basic.uniformDegreeBound}
  \uses{def:ambient_setup}
  % SOURCE: Section 6 of \cite{MR4276287-mordell-lang-curves} (read from references/MR4276287-mordell-lang-curves-source/SettingUp.tex)
  % SOURCE QUOTE: "There exists a constant \(c\) such that ... \(\deg(\mathfrak{C}_s-P) \le c\) ... and \(h(\mathfrak{C}_s-P_s) \le c \max\{1, h_{\overline{\mathbb{A}_g},\cM}(\tau(s))\}\)."
  % SOURCE QUOTE PROOF: The section bounds both the projective degree and the moduli height of the translated fibers uniformly in the base point.
  \textit{Source: Section 6 of \cite{MR4276287-mordell-lang-curves} (read from references/MR4276287-mordell-lang-curves-source/SettingUp.tex).}
  Translates of a fixed fiber curve have uniformly bounded projective degree, and
  one translate in each fiber has height controlled by the moduli point.
\end{lemma}

% SOURCE QUOTE PROOF: "We need a quasi-section of $\mathfrak{C}_g\rightarrow\mathbb{M}_g$ as provided by \cite[Corollaire~17.16.3(ii)]{EGAIV}. So there is an affine scheme $S$ and a morphism $S\rightarrow \mathfrak{C}_g$ that factors through a surjective, quasi-finite, \'etale morphism $S\rightarrow \mathbb{M}_g$. We consider the product $\mathfrak{C}_g\times_{\mathbb{M}_g} S\rightarrow \mathfrak{C}_g^{[2]}$ composed with the Faltings--Zhang morphism $\mathscr{D}_1\colon \mathfrak{C}_g^{[2]}\rightarrow\mathfrak{A}_g\times_{\IA_g} \mathbb{M}_g$ over $\mathbb{M}_g$ and then the projection of $\mathfrak{A}_g$." (read from references/MR4276287-mordell-lang-curves-source/SettingUp.tex)
\begin{proof}
  \uses{def:ambient_setup}
  By \cite[Corollaire~17.16.3(ii)]{EGAIV}, there exists a quasi-section of
  \(\mathfrak{C}_g \to \mathbb{M}_g\): an affine scheme \(S_0\) with a morphism
  \(S_0 \to \mathfrak{C}_g\) factoring through a surjective quasi-finite étale
  \(S_0 \to \mathbb{M}_g\).
  The composition \(\mathfrak{C}_g \times_{\mathbb{M}_g} S_0 \to \mathfrak{A}_g\)
  via the Faltings--Zhang morphism \(\mathscr{D}_1\) has image a constructible
  union of finitely many irreducible locally closed subsets \(X_i \subset \mathfrak{A}_g\).
  For any \(s \in \mathbb{M}_g(\overline{\mathbb{Q}})\), the curve \(\mathfrak{C}_s - P_s\)
  arises as an irreducible component of the intersection of some \(\overline{X_i}\)
  with the fiber \(\mathbb{P}^n \times \{\tau(s)\}\).
  Part (i) follows from Bézout's theorem applied to the Segre embedding
  \(\mathbb{P}^n \times \mathbb{P}^m \to \mathbb{P}^{(n+1)(m+1)-1}\):
  the degree of \(\mathfrak{C}_s - P_s\) is bounded above uniformly in \(s\),
  and translating by any \(P \in \mathfrak{A}_{g,\tau(s)}(\overline{\mathbb{Q}})\)
  does not change the degree.
  Part (ii) follows from the Arithmetic Bézout theorem \cite[Théorème~3]{HauteursAlt3}:
  the height of \(\overline{X_i}\) and the degree of \(\mathbb{P}^n \times \{\tau(s)\}\)
  are bounded independently of \(s\), while the height of the fiber hyperplane is
  bounded linearly in \(h(\tau(s))\), hence linearly in
  \(h_{\overline{\mathbb{A}_g},\mathcal{M}}(\tau(s))\) by the Height Machine.
\end{proof}

\subsection{Non-degeneracy of $\mathscr{D}_M(\mathfrak{C}_S^{[M+1]})$ for large $M$}

\begin{theorem}[Non-degeneracy of the Faltings--Zhang image]
  \label{thm:nondeg_for_bd}
  \lean{MR4276287UniformityInMordellLangForCurves.Basic.faltingsZhangNondegenerate}
  \uses{def:ambient_setup,def:nondegenerate_app,lem:curve_generates_jacobian,def:faltings_zhang}
  % SOURCE: Section 6 of \cite{MR4276287-mordell-lang-curves} (read from references/MR4276287-mordell-lang-curves-source/SettingUp.tex)
  % SOURCE QUOTE: "Let \(S\) be an irreducible variety with a (not necessarily dominant) quasi-finite morphism \(S \rightarrow \mathbb{M}_g\). Assume \(g \ge 2\) and \(M \ge 3g-2\). Then \(\mathscr{D}_M( \mathfrak{C}_S^{[M+1]} )\) ... is non-degenerate."
  % SOURCE QUOTE PROOF: Gao's theorem gives the needed non-degeneracy once the fiber power is large enough.
  \textit{Source: Section 6 of \cite{MR4276287-mordell-lang-curves} (read from references/MR4276287-mordell-lang-curves-source/SettingUp.tex).}
  For sufficiently large fiber power, Gao's theorem shows that the Faltings--Zhang
  image of the universal curve family is non-degenerate.
\end{theorem}

% SOURCE QUOTE PROOF: "This theorem, essentially \cite[Theorem~1.2']{GaoBettiRank}, is a consequence of Theorem~1.3 of \textit{loc.cit.} Because of its importance to the current paper, we hereby give more details of the deduction. We start by showing the result for the case where $\mathfrak{C}_S \rightarrow S$ admits a section $\epsilon$. In this case $\epsilon$ induces an Abel--Jacobi embedding $j_{\epsilon} \colon \mathfrak{C}_S \rightarrow \mathrm{Jac}(\mathfrak{C}_S/S)$, which is a closed immersion of $S$-schemes." (read from references/MR4276287-mordell-lang-curves-source/SettingUp.tex)
\begin{proof}
  \uses{def:ambient_setup, def:nondegenerate_app, prop:betti_map}
  The theorem is essentially \cite[Theorem~1.2']{GaoBettiRank}, deduced from
  \cite[Theorem~1.3]{GaoBettiRank}.
  First assume \(\mathfrak{C}_S \to S\) admits a section \(\epsilon\).
  Then \(\epsilon\) induces an Abel--Jacobi embedding
  \(j_\epsilon \colon \mathfrak{C}_S \to \mathrm{Jac}(\mathfrak{C}_S/S)\),
  and the modular map \(\iota \colon \mathrm{Jac}(\mathfrak{C}_S/S) \to \mathfrak{A}_g\)
  is quasi-finite (since the Torelli map \(\tau\) is quasi-finite by
  \cite[Lemma~1.11]{OortSteenbrink}).
  Applying \cite[Theorem~1.3(ii)]{GaoBettiRank} to \(j_\epsilon(\mathfrak{C}_S)\)
  requires: (a) its image under \(\iota\) is generically finite (holds since \(\iota\) is quasi-finite);
  (b) each fiber curve generates its Jacobian (true for smooth curves); (c) \(g \ge 2\).
  Under these hypotheses, \(\mathscr{D}_M(\mathfrak{C}_S^{[M+1]})\) is non-degenerate
  for \(M \ge \dim S + 1 \le 3g - 2\).
  For general \(S\) without a section, pass to a quasi-finite étale dominant cover
  \(\rho \colon S' \to S\) on which \(\mathfrak{C}_{S'} = \mathfrak{C}_S \times_S S'\)
  admits a section. The non-degeneracy of
  \(X' = \mathscr{D}_M(\mathfrak{C}_{S'}^{[M+1]})\)
  in \(\mathcal{A}' = \mathfrak{A}_g^{[M]} \times_{\mathbb{A}_g} S'\) transfers
  to \(X = \mathscr{D}_M(\mathfrak{C}_S^{[M+1]})\) in \(\mathcal{A}\)
  via the bianalytic identification of \(\Delta' \cong \Delta\) and uniqueness of
  Betti maps.
\end{proof}

\subsection{Technical lemmas}

\begin{lemma}[N-Alon combinatorial lemma]
  \label{lem:NAlon}
  \lean{MR4276287UniformityInMordellLangForCurves.Basic.nAlon}
  \uses{def:ambient_setup}
  % SOURCE: Section 6 of \cite{MR4276287-mordell-lang-curves} (read from references/MR4276287-mordell-lang-curves-source/SettingUp.tex)
  % SOURCE QUOTE: "Let \(C \subset \IP_k^n\) be an irreducible curve ... Then there exists a number \(B\), depending only on \(M, \deg C,\) and \(\deg Z\), ... if \(\Sigma \subset C(k)\) has cardinality \(\ge B\), then \(\Sigma^M \not\subseteq Z(k)\)."
  % SOURCE QUOTE PROOF: This is the product-avoidance input used to keep large subsets of a curve from collapsing into a proper closed set.
  \textit{Source: Section 6 of \cite{MR4276287-mordell-lang-curves} (read from references/MR4276287-mordell-lang-curves-source/SettingUp.tex).}
  Large enough subsets of a curve cannot have all \(M\)-fold products trapped in
  a fixed proper closed subset of the ambient product.
\end{lemma}

% SOURCE QUOTE PROOF: "Let us prove this lemma by induction on $M$. The case $M = 1$ follows easily from B\'{e}zout's Theorem. Assume the lemma is proved for $1, \ldots, M-1$. Let $q \colon (\IP_k^n)^M \rightarrow \IP_k^n$ be the projection to the first factor. The number of irreducible components of $Z \cap C^M$ and their degrees are bounded from above in terms of $M,\deg C,$ and $\deg Z$ by B\'ezout's Theorem applied to the Segre embedding." (read from references/MR4276287-mordell-lang-curves-source/SettingUp.tex)
\begin{proof}
  \uses{def:ambient_setup}
  The proof is by induction on \(M\). The base case \(M = 1\) follows from Bézout's theorem:
  the finite set \(Z \cap C\) has degree bounded by \(\deg Z \cdot \deg C\), so any
  \(\Sigma \subset C(k)\) with \(|\Sigma| > \deg Z \cdot \deg C\) satisfies
  \(\Sigma \not\subset Z(k)\).
  For the inductive step, let \(q \colon (\mathbb{P}^n)^M \to \mathbb{P}^n\)
  be the projection to the first factor.
  Decompose \(Z \cap C^M\) into components with \(\dim q(\text{component}) \ge 1\)
  (forming \(Z'\)) and components where \(q\) maps to a point (forming \(Z''\)).
  The number of components in \(Z''\) is bounded by Bézout, so \(q(Z'')\) is a finite
  set of cardinality \(\le B''\).
  For each \(P \in C(k)\), the fiber \(q|_{Z'}^{-1}(P) \cap C^{M-1}\) does not contain
  \(C^{M-1}\) (since \(\dim q|_{Z'}^{-1}(P) \le \dim Z' - 1 \le M - 2\)),
  so by the induction hypothesis there exists \(B'\) such that any
  \(\Sigma\) of cardinality \(\ge B'\) satisfies
  \(\{P\} \times \Sigma^{M-1} \not\subset Z'(k)\) for every \(P \in C(k)\).
  Taking \(B = \max\{B', B'' + 1\}\) completes the induction.
\end{proof}

\begin{lemma}[Packets alternative]
  \label{lem:packets_alternative3}
  \lean{MR4276287UniformityInMordellLangForCurves.Basic.packetsAlternative}
  \uses{def:ambient_setup}
  % SOURCE: Section 6 of \cite{MR4276287-mordell-lang-curves} (read from references/MR4276287-mordell-lang-curves-source/SettingUp.tex)
  % SOURCE QUOTE: "If \(Z\) is an irreducible Zariski closed and proper subset of \(\mathscr{D}_M(C^{M+1})\), then \(\#\{ P \in C(\overline{\IQ}) : (C-P)^M \subset Z \} \le 84(g-1)\)."
  % SOURCE QUOTE PROOF: Hurwitz theory bounds the number of points whose packet condition lands in a fixed proper closed subset.
  \textit{Source: Section 6 of \cite{MR4276287-mordell-lang-curves} (read from references/MR4276287-mordell-lang-curves-source/SettingUp.tex).}
  A proper closed condition on the Faltings--Zhang image can only hold for
  uniformly many base points on a curve of genus at least \(2\).
\end{lemma}

% SOURCE QUOTE PROOF: "For simplicity denote by $\Xi = \{ P \in C(\overline{\IQ}) : (C-P)^M \subset Z \}$. Fix $P_0\in \Xi$. It suffices to prove that there are only $84(g-1)$ possibilities for $P_1 - P_0$ when $P_1$ runs over $\Xi$. Say $P_1 \in \Xi$ and let $i\in \{0,1\}$. Note that $Z\subsetneq \mathscr{D}_M(C^{M+1})$ and so $\dim Z < \dim \mathscr{D}_M(C^{M+1})\le M+1$, as $\mathscr{D}_M(C^{M+1})$ irreducible. Note that $(C-P_i)^M\subset Z$, and so both have dimension $M$. As $Z$ is irreducible we find $(C - P_i)^M=Z$ for $i=0$ and $i=1$. Applying the first projection $A^M\rightarrow A$ yields $C-P_1=C-P_0$. In other words, $P_1-P_0$ stabilizes $C$. By Hurwitz's Theorem \cite{Hurwitz1892}, a smooth curve over $\IQbar$ of genus $g\ge 2$ has at most $84(g-1)$ automorphisms." (read from references/MR4276287-mordell-lang-curves-source/SettingUp.tex)
\begin{proof}
  \uses{def:ambient_setup}
  Let \(\Xi = \{P \in C(\overline{\mathbb{Q}}) : (C-P)^M \subset Z\}\).
  Fix \(P_0 \in \Xi\); it suffices to bound the number of possibilities for \(P_1 - P_0\)
  as \(P_1\) ranges over \(\Xi\).
  For any \(P_1 \in \Xi\) and \(i \in \{0,1\}\), since
  \(Z \subsetneq \mathscr{D}_M(C^{M+1})\) and \(\dim Z < M+1\),
  while \((C-P_i)^M \subset Z\) has dimension \(M\), irreducibility of \(Z\)
  forces \((C - P_i)^M = Z\) for both \(i = 0\) and \(i = 1\).
  Projecting onto the first factor \(A^M \to A\) gives \(C - P_1 = C - P_0\),
  hence \(P_1 - P_0\) is an automorphism of \(C\).
  By Hurwitz's theorem \cite{Hurwitz1892}, a smooth curve of genus \(g \ge 2\)
  defined over \(\overline{\mathbb{Q}}\) has at most \(84(g-1)\) automorphisms.
  Therefore \(\#\Xi \le 84(g-1)\).
\end{proof}

\section{N\'{e}ron--Tate distance between points on curves}

\begin{proposition}[Points on curves are sparse]
  \label{prop:alg_pt_far}
  \lean{MR4276287UniformityInMordellLangForCurves.Basic.algPointsFar}
  \uses{thm:nondeg_for_bd,thm:ht_inequality_full,lem:uniform_degree_bound,lem:NAlon,lem:packets_alternative3}
  % SOURCE: Section 7 of \cite{MR4276287-mordell-lang-curves} (read from references/MR4276287-mordell-lang-curves-source/DistanceCurve.tex)
  % SOURCE QUOTE: "There exist positive constants \(c_1, c_2, c_3, c_4\) ... either \(P \in \Xi_s\), or \(\#\{Q \in \mathfrak{C}_s(\IQbar) : \hat{h}(Q-P) \le h_{\overline{\mathbb{A}_g},\cM}(\tau(s))/c_3\} < c_4\)."
  % SOURCE QUOTE PROOF: This is the distance alternative obtained by combining the non-degeneracy input, the height inequalities, and the counting lemmas.
  \textit{Source: Section 7 of \cite{MR4276287-mordell-lang-curves} (read from references/MR4276287-mordell-lang-curves-source/DistanceCurve.tex).}
  For a family of genus-\(g\) curves with large moduli height, every point is
  either exceptional or has only uniformly boundedly many nearby points in the
  canonical metric.
\end{proposition}

% SOURCE QUOTE PROOF: "We fix an immersion of $\mathbb{M}_g$ into some projective space and let $\overline{\mathbb{M}_g}$ denote its Zariski closure. By a standard triangle inequality estimate, there exist constants $c''>0$ and $c'''\ge 0$ such that $h(s) \ge c'' h_{\overline{\mathbb{A}_g},\cM}(\tau(s)) - c'''$ for all $s \in \mathbb{M}_g(\IQbar)$; see \cite[Lemma~4]{Silverman:heightest11}. Our proof is by induction on $\dim S$." (read from references/MR4276287-mordell-lang-curves-source/DistanceCurve.tex)
\begin{proof}
  \uses{thm:nondeg_for_bd, thm:ht_inequality_full, lem:uniform_degree_bound, lem:NAlon, lem:packets_alternative3}
  The proof is by induction on \(\dim S\).
  The base case \(\dim S = 0\) is handled by enlarging \(c_1\).
  For \(\dim S \ge 1\), fix \(M = 3g - 2\). By \cref{thm:nondeg_for_bd},
  the Zariski closed subvariety
  \(X = \mathscr{D}_M(\mathfrak{C}_{S^{\mathrm{sm}}}^{[M+1]})\)
  of \(\mathcal{A} = \mathfrak{A}_g^{[M]} \times_{\mathbb{A}_g} S^{\mathrm{sm}} \to S^{\mathrm{sm}}\)
  is non-degenerate. Applying \cref{thm:ht_inequality_full} to \(\mathcal{A}\) and \(X\)
  gives a Zariski open dense \(U \subset X\) and constants \(c, c' > 0\) such that,
  for \(s\) with \(h_{\overline{\mathbb{A}_g},\mathcal{M}}(\tau(s)) \ge c_1\)
  and points \((Q_1 - P, \ldots, Q_M - P) \in U(\overline{\mathbb{Q}})\),
  one has \(c\,h(s) \le \hat{h}(Q_1 - P) + \cdots + \hat{h}(Q_M - P) + c'\).
  The complement \(W = X \setminus U\) is a proper closed subset.
  For each irreducible component \(Z\) of \(W_s\),
  \cref{lem:packets_alternative3} gives
  \(\#\Xi_Z = \#\{P \in \mathfrak{C}_s(\overline{\mathbb{Q}}) : (\mathfrak{C}_s - P)^M \subset Z\} \le 84(g-1)\).
  Setting \(\Xi_s = \bigcup_Z \Xi_Z\) gives \(\#\Xi_s \le c_2\) for a constant
  \(c_2\) independent of \(s\).
  For \(P \notin \Xi_s\), one has \((\mathfrak{C}_s - P)^M \not\subset W_s\).
  By \cref{lem:uniform_degree_bound} and \cref{lem:NAlon}, there is a uniform
  bound \(c_4\) such that any \(\Sigma \subset \mathfrak{C}_s(\overline{\mathbb{Q}})\)
  with \(|\Sigma| \ge c_4\) satisfies \((\Sigma - P)^M \not\subset W_s\).
  Taking \(\Sigma = \{Q : \hat{h}(Q - P) \le h_{\overline{\mathbb{A}_g},\mathcal{M}}(\tau(s))/c_3\}\)
  with \(c_3 = 2M/c\), if \(\#\Sigma \ge c_4\) one extracts
  \((Q_1 - P, \ldots, Q_M - P) \in U(\overline{\mathbb{Q}})\),
  applies the height inequality to bound \(h(s)\) above by \(2c'/c\),
  and concludes \(h_{\overline{\mathbb{A}_g},\mathcal{M}}(\tau(s)) < c_1\)
  for \(c_1\) large enough. This contradicts the hypothesis,
  so \(\#\Sigma < c_4\), giving alternative (ii).
  The induction hypothesis handles \(s\) above the proper closed
  subvarieties \(S_1, \ldots, S_r\) comprising the Zariski closure of \(S \setminus \pi(U)\).
\end{proof}

\section{Proof of Theorems~\ref{ThmBdRatIntro}, \ref{ThmBdFinRank}, and \ref{thm:BdTorIntroNF}}

\begin{lemma}[Quantitative Vojta--Mumford estimate]
  \label{lem:vojtamumford}
  \lean{MR4276287UniformityInMordellLangForCurves.Basic.vojtaMumford}
  \uses{thm:silverman_tate,lem:uniform_degree_bound}
  % SOURCE: Section 8 of \cite{MR4276287-mordell-lang-curves} (read from references/MR4276287-mordell-lang-curves-source/RatPt.tex)
  % SOURCE QUOTE: "Let \(C\) be an irreducible curve in \(A\). There exists a constant \(c=c(n,\deg C)\ge 1\) ... \# \{P \in C(\IQbar)\cap \Gamma : \hat h_{L}(P) > c\max\{1,h(C),c_{\mathrm{NT}},h_1\} \} \le c^\rho."
  % SOURCE QUOTE PROOF: The estimate is the large-points part of the Vojta--Mumford method and depends on the subgroup rank through the exponent.
  \textit{Source: Section 8 of \cite{MR4276287-mordell-lang-curves} (read from references/MR4276287-mordell-lang-curves-source/RatPt.tex).}
  A curve inside an abelian variety contributes at most exponentially many points
  of large canonical height, with the exponent controlled by the subgroup rank.
\end{lemma}

% SOURCE QUOTE PROOF: "The lemma below is \cite[Corollary~2.3]{DGH1p} which is itself a standard application of R\'emond's explicit formulation of the Vojta and Mumford inequalities. We thus obtain a bound that is exponential in the rank of the subgroup $\Gamma$ for points of sufficiently large N\'eron--Tate height." (read from references/MR4276287-mordell-lang-curves-source/RatPt.tex)
\begin{proof}
  \uses{thm:silverman_tate, lem:uniform_degree_bound}
  This is \cite[Corollary~2.3]{DGH1p}, a standard application of Rémond's
  explicit formulation \cite{Remond:Decompte, remond:vojtasup} of the
  Vojta and Mumford inequalities.
  The Vojta inequality provides, for each point \(P \in C(\overline{\mathbb{Q}}) \cap \Gamma\)
  of sufficiently large Néron--Tate height, a strong linear constraint on
  the position of \(P\) relative to other points in \(\Gamma\).
  The Mumford inequality gives a separation bound: two distinct such high-height points
  are far apart in the Néron--Tate metric.
  Combining these two inequalities shows that there can be at most \(c^\rho\)
  points of \(C(\overline{\mathbb{Q}}) \cap \Gamma\) with
  \(\hat{h}_L(P) > c\max\{1, h(C), c_{\mathrm{NT}}, h_1\}\),
  where \(c\) depends only on \(n\) and \(\deg C\), and \(h_1\) is a height
  measuring the complexity of the group law on \(A\) in projective coordinates.
  The \cref{thm:silverman_tate} controls \(c_{\mathrm{NT}}\) in terms of
  \(h_{\overline{S}}(\pi(P))\), and \cref{lem:uniform_degree_bound} controls
  \(h(C)\) in terms of the moduli height.
\end{proof}

\section{The Silverman--Tate Theorem revisited}

\begin{lemma}[Silverman--Tate height comparison]
  \label{thm:silverman_tate}
  \lean{MR4276287UniformityInMordellLangForCurves.Basic.silvermanTate}
  \uses{def:ambient_setup,def:nt_height,def:height_fib}
  % SOURCE: Appendix A of \cite{MR4276287-mordell-lang-curves} (read from references/MR4276287-mordell-lang-curves-source/silvermantate.tex)
  % SOURCE QUOTE: "There exists a constant \(c>0\) such that for all \(P \in \cA(\IQbar)\) we have \(\bigl| \hat h_{\cA}(P) - \hcA(P)\bigr| \le c \max\{1,h_{\overline S}(\pi(P))\}.\)"
  % SOURCE QUOTE PROOF: The comparison is the canonical-height versus projective-height estimate used to move between the fiber and the base.
  \textit{Source: Appendix A of \cite{MR4276287-mordell-lang-curves} (read from references/MR4276287-mordell-lang-curves-source/silvermantate.tex).}
  The canonical height on a fiber differs from the naive projective height by an
  error bounded linearly in the height of the base point.
\end{lemma}

% SOURCE QUOTE PROOF: "Having (\ref{eq:duplicationbound}) at our disposal the proof follows a well-known argument. Indeed, say $l\ge k\ge 0$ are integers. Then applying the triangle inequality to the appropriate telescoping sum yields $\left|\frac{\hcA([2^l](P))}{4^l}- \frac{\hcA([2^k](P))}{4^k}\right| \le \sum_{m=k}^{l-1} 4^{-(m+1)} \left|{\hcA([2^{m+1}](P))}- 4{\hcA([2^m](P))}\right|$. We apply (\ref{eq:duplicationbound}) to $[2^m](P)$ and find that the sum is bounded by $c_1 x \sum_{m=k}^{l-1}4^{-(m+1)} \le c_1 x 4^{-k}$ where $x = \max\{1,h_{\overline S}(\pi(P))\}$. So $\left(\hcA([2^l](P))/4^l\right)_{l\ge 1}$ is a Cauchy sequence with limit $\hat h_{\overline\cA}(P)$. Taking $k=0$ and $l\rightarrow \infty$ we obtain from the estimates above that $|\hat h_{\overline\cA}(P) - \hcA(P)|\le c_1x$, as desired." (read from references/MR4276287-mordell-lang-curves-source/silvermantate.tex)
\begin{proof}
  \uses{def:ambient_setup}
  The key input is the duplication bound: there exists \(c_1 > 0\) such that
  \(\bigl|\hcA([2](P)) - 4\hcA(P)\bigr| \le c_1 \max\{1, h_{\overline{S}}(\pi(P))\}\)
  for all \(P \in \mathcal{A}(\overline{\mathbb{Q}})\).
  This is proved by constructing the graph of \([2]\) inside
  \(\overline{\mathcal{A}} \times_{\overline{S}} \overline{\mathcal{A}}\),
  extending the difference line bundle \([2]^*\overline{\mathcal{L}} \otimes \rho^*\overline{\mathcal{L}}^{\otimes(-4)}\)
  to the Zariski closure, and applying the Height Machine after desingularizing
  the base via Hironaka's theorem.
  Given the duplication bound, the Néron--Tate height is defined as
  \(\hat{h}_{\mathcal{A}}(P) = \lim_{N \to \infty} \hcA([2^N](P))/4^N\).
  The telescoping estimate
  \(\bigl|\hcA([2^l](P))/4^l - \hcA([2^k](P))/4^k\bigr| \le c_1 \max\{1,h_{\overline{S}}(\pi(P))\} \cdot 4^{-k}\)
  shows the sequence is Cauchy; taking \(k = 0\) and \(l \to \infty\) gives
  \(\bigl|\hat{h}_{\mathcal{A}}(P) - \hcA(P)\bigr| \le c_1 \max\{1, h_{\overline{S}}(\pi(P))\}\).
\end{proof}

\section{Full version of Theorem~\ref{ThmHtInequality}}

\subsection{Betti map}

\begin{proposition}[Betti map without the level-structure restrictions]
  \label{prop:betti_map_app}
  \lean{MR4276287UniformityInMordellLangForCurves.Basic.bettiMapApp}
  \uses{def:ambient_setup}
  % SOURCE: Appendix B of \cite{MR4276287-mordell-lang-curves} (read from references/MR4276287-mordell-lang-curves-source/HtIneqFinalVer.tex)
  % SOURCE QUOTE: "Let \(s_0 \in S(\IC)\). Then there exist an open neighborhood \(\Delta\) of \(s_0\) in \(S^{\mathrm{an}}\), and a map \(b_{\Delta} \colon \cA_{\Delta} := \pi^{-1}(\Delta) \rightarrow \mathbb T^{2g}\), called the Betti map, with the following properties."
  % SOURCE QUOTE PROOF: The appendix repeats the Betti-map construction after passing to a finite étale cover when needed.
  \textit{Source: Appendix B of \cite{MR4276287-mordell-lang-curves} (read from references/MR4276287-mordell-lang-curves-source/HtIneqFinalVer.tex).}
  The local Betti-map construction persists for arbitrary regular bases and
  abelian schemes after passing to a finite étale cover when necessary.
\end{proposition}

% SOURCE QUOTE PROOF: "Our proof of Proposition~\ref{PropBettiMapApp} follows the construction in \cite[$\mathsection$3-$\mathsection$4]{GaoBettiRank}. We divide it into several steps. By \cite[$\mathsection$2.1]{GenestierNgo}, $\cA \rightarrow S$ carries a polarization of type $D = \mathrm{diag}(d_1,\ldots,d_g)$ for some positive integers $d_1|d_2|\cdots|d_g$. \textbf{Case: Moduli space with level structure.} Fix $\ell \ge \bestlevel$ with $(\ell,d_g) = 1$. We start by proving Proposition~\ref{PropBettiMapApp} for $S = \mathbb A_{g,D,\ell}$, the moduli space of abelian varieties of dimension $g$ polarized of type $D$ with level-$\ell$-structure." (read from references/MR4276287-mordell-lang-curves-source/HtIneqFinalVer.tex)
\begin{proof}
  \uses{def:ambient_setup, prop:betti_map}
  The proof proceeds in three cases.
  \textit{Case 1: Fine moduli space with level structure.}
  Fix \(\ell \ge 3\) with \((\ell, d_g) = 1\). When
  \(S = \mathbb{A}_{g,D,\ell}\) is the fine moduli space of abelian varieties
  polarized of type \(D = \mathrm{diag}(d_1,\ldots,d_g)\) with level-\(\ell\)-structure,
  the same Siegel upper half-space construction as in \cref{prop:betti_map} applies:
  the universal covering \(\mathfrak{H}_g \to \mathbb{A}_{g,D,\ell}^{\mathrm{an}}\)
  gives the Betti map via the inverse of \((a,b,Z) \mapsto (Da + Zb, Z)\).
  \textit{Case 2: Base carries level structure.}
  If \(\mathcal{A} \to S\) already carries level-\(\ell\)-structure, the fine
  moduli property gives a Cartesian diagram \(\mathcal{A} \to \mathfrak{A}_{g,D,\ell}\)
  over \(S \to \mathbb{A}_{g,D,\ell}\); the Betti map is pulled back along this diagram.
  \textit{Case 3: General case.}
  Fix a prime \(\ell \ge 3\) with \((\ell, d_g) = 1\). Take an irreducible component
  \(S_0\) of \(\ker[\ell] \subset \mathcal{A}\); since \(S\) is regular,
  \(S_0 \to S\) is finite and étale. Repeating finitely many times yields a
  finite étale cover \(\rho \colon S' \to S\) with \(S'\) irreducible and regular,
  and \(\mathcal{A}' = \mathcal{A} \times_S S' \to S'\) carrying level-\(\ell\)-structure.
  Apply Case 2 to get a Betti map \(b_{\Delta'}\) on \(\mathcal{A}'_{\Delta'}\).
  Shrink \(\Delta'\) so that \(\rho|_{\Delta'}\) is a homeomorphism onto
  \(\Delta = \rho(\Delta')\); then \(b_\Delta := b_{\Delta'} \circ (\mathcal{A}'_{\Delta'} \cong \mathcal{A}_\Delta)^{-1}\)
  is the desired Betti map.
\end{proof}

\begin{lemma}[Betti rank is insensitive to shrinking the base]
  \label{lem:betti_rank_zar_open_app}
  \lean{MR4276287UniformityInMordellLangForCurves.Basic.bettiRankZarOpen}
  \uses{prop:betti_map_app}
  % SOURCE: Appendix B of \cite{MR4276287-mordell-lang-curves} (read from references/MR4276287-mordell-lang-curves-source/HtIneqFinalVer.tex)
  % SOURCE QUOTE: "Let \(b_{\Delta} \colon \cA_{\Delta} \rightarrow \mathbb T^{2g}\) be a Betti map as in Proposition~\ref{PropBettiMapApp}. Let \(X\) be an irreducible subvariety of \(\cA\) with \(\an{X}\cap\cA_\Delta\not=\emptyset\). Let \(U\) be a Zariski open dense subset of \(X\). Then ... the maximum rank is the same on \(X\) and on \(U\)."
  % SOURCE QUOTE PROOF: The same local rank is attained on a dense open subset because rank is locally constant where it is maximal.
  \textit{Source: Appendix B of \cite{MR4276287-mordell-lang-curves} (read from references/MR4276287-mordell-lang-curves-source/HtIneqFinalVer.tex).}
  Passing to a Zariski open dense subset does not change the maximal Betti-map
  rank seen on the smooth locus.
\end{lemma}

% SOURCE QUOTE PROOF: "The statement \eqref{EqBettiRankSmallerUnderRestriction} is true on replacing ``$=$'' by ``$\ge$'', as $\sman{X}\supset \sman{U}$. For the converse inequality we set $\max_{x \in \sman{X} \cap \cA_{\Delta}} \mathrm{rank}_{\IR} (\mathrm{d}b_{\Delta}|_{X^{\mathrm{sm,an}}})_x = r$ and pick $x \in \sman{X} \cap \cA_{\Delta}$ satisfying $\mathrm{rank}_{\IR} (\mathrm{d}b_{\Delta}|_{X^{\mathrm{sm,an}}})_x = r$. Then there exists an open neighborhood $V$ of $x$ in $\sman{X}$ such that $\mathrm{rank}_{\IR} (\mathrm{d}b_{\Delta}|_{X^{\mathrm{sm,an}}})_u = r$ for all $u \in V$. But $U^{\mathrm{sm}}(\IC) \cap V \not= \emptyset$ since $U^{\mathrm{sm}}\not=\emptyset$ is Zariski open in $X$ and $V$ is Zariski dense in $X$." (read from references/MR4276287-mordell-lang-curves-source/HtIneqFinalVer.tex)
\begin{proof}
  \uses{prop:betti_map_app}
  The inequality \(\ge\) is immediate since \(X^{\mathrm{sm}} \supset U^{\mathrm{sm}}\).
  For the reverse, let \(r\) be the maximum rank on \(X^{\mathrm{sm}} \cap \mathcal{A}_\Delta\)
  and pick \(x \in X^{\mathrm{sm}} \cap \mathcal{A}_\Delta\) where rank \(r\) is attained.
  Since the rank function is real-analytic and \(b_\Delta\) is real-analytic, the rank
  equals \(r\) on an open neighborhood \(V\) of \(x\) in \(X^{\mathrm{sm}}\).
  Because \(U^{\mathrm{sm}}\) is Zariski open and dense in \(X\), it is analytically dense
  in \(X^{\mathrm{an}}\), so \(U^{\mathrm{sm}}(\mathbb{C}) \cap V \ne \emptyset\).
  Any \(u \in U^{\mathrm{sm}}(\mathbb{C}) \cap V\) achieves rank \(r\)
  for \(b_\Delta|_{U^{\mathrm{sm,an}}}\), giving the reverse inequality.
\end{proof}

\subsection{Non-degenerate subvariety and Theorem~\ref{ThmHtInequality}}

\begin{definition}[Non-degenerate subvariety]
  \label{def:nondegenerate_app}
  \lean{MR4276287UniformityInMordellLangForCurves.Basic.nonDegenerate}
  \uses{prop:betti_map_app}
  % SOURCE: Appendix B of \cite{MR4276287-mordell-lang-curves} (read from references/MR4276287-mordell-lang-curves-source/HtIneqFinalVer.tex)
  % SOURCE QUOTE: "An irreducible subvariety \(X\) of \(\cA\) is said to be non-degenerate if there exists an open non-empty subset \(\Delta\) ... such that ... \(\max_{x \in X^{\mathrm{sm}}(\IC) \cap \cA_{\Delta}} \mathrm{rank}_{\IR} (\mathrm{d}b_{\Delta}|_{X^{\mathrm{sm,an}}})_x = 2\dim X\)."
  % SOURCE QUOTE PROOF: This restates the appendix criterion in the language of local Betti rank on the smooth locus.
  \textit{Source: Appendix B of \cite{MR4276287-mordell-lang-curves} (read from references/MR4276287-mordell-lang-curves-source/HtIneqFinalVer.tex).}
  An irreducible subvariety is non-degenerate if some local Betti map has full
  real rank \(2\dim X\) on a smooth point of the analytification.
\end{definition}

\begin{theorem}[Full height inequality]
  \label{thm:ht_inequality_full}
  \lean{MR4276287UniformityInMordellLangForCurves.Basic.heightInequalityFull}
  \uses{prop:betti_map_app,lem:betti_rank_zar_open_app,def:nondegenerate_app,thm:ht_inequality}
  % SOURCE: Appendix B of \cite{MR4276287-mordell-lang-curves} (read from references/MR4276287-mordell-lang-curves-source/HtIneqFinalVer.tex)
  % SOURCE QUOTE: "Let \(X\) be an irreducible subvariety of \(\cA\) ... Then there exist constants \(c_1>0\) and \(c_2\ge 0\) and a Zariski open dense subset \(U\) of \(X\) with \(\hat{h}_{\cA,\cL}(P) \ge c_1 h_{\overline{S},\cM}( \pi(P) ) - c_2\)."
  % SOURCE QUOTE PROOF: The appendix removes the extra Hyp assumptions by passing through the general Betti-map setup and then shrinking to a dense open subset.
  \textit{Source: Appendix B of \cite{MR4276287-mordell-lang-curves} (read from references/MR4276287-mordell-lang-curves-source/HtIneqFinalVer.tex).}
  The appendix first enlarges the Betti-map input to arbitrary regular bases and
  then performs a six-step devissage: close up \(X\), replace the base by the
  image of \(X\), pass to a regular cover carrying level structure and a
  principal polarization, and finally invoke the height inequality already
  proved under the paper's hypotheses. The result is the same lower bound for
  the canonical height on a dense open subset of any non-degenerate subvariety.
\end{theorem}

% SOURCE QUOTE PROOF: "We will reduce the current theorem to Theorem~\ref{ThmHtInequality} by successively assuming, in addition to the hypothesis of Theorem~\ref{ThmHtInequalityFull}, that (i) $X$ is Zariski closed in $\cA$, (ii) $\pi|_X \colon X\rightarrow S$ is dominant, (iii) $S$ is regular, (iv) $\cA \rightarrow S$ is $S$-isogenous to an abelian scheme which carries a principal polarization, (v) $\cA \rightarrow S$ carries a principal polarization, (vi) $\cA$ carries a level $\ell$-structure, and (vii) we have the same hypothesis as Theorem~\ref{ThmHtInequality}. We will proceed the proof with six d\'evissage steps." (read from references/MR4276287-mordell-lang-curves-source/HtIneqFinalVer.tex)
\begin{proof}
  \uses{prop:betti_map_app, lem:betti_rank_zar_open_app, def:nondegenerate_app, thm:ht_inequality}
  The proof reduces to \cref{thm:ht_inequality} via six devissage steps.
  \textit{Step (i):} Replace \(X\) by its Zariski closure \(\overline{X}\) in \(\mathcal{A}\);
  non-degeneracy and the conclusion are preserved.
  \textit{Step (ii):} Restrict to \(S' = \pi(X)\) with reduced structure; \(X\) dominates \(S'\)
  and non-degeneracy pulls back unchanged.
  \textit{Step (iii):} Replace \(S\) by its regular locus \(S^{\mathrm{sm}}\) and \(X\) by
  \(X \cap \pi^{-1}(S^{\mathrm{sm}})\), which is still non-degenerate by \cref{lem:betti_rank_zar_open_app}.
  \textit{Step (iv):} By \cite[Section~23, Corollary~1]{MumfordAbVar}, pass to a
  quasi-finite étale dominant cover \(\rho \colon S' \to S\) such that the generic fiber of
  \(\mathcal{A}' = \mathcal{A} \times_S S'\) is isogenous to a principally polarized abelian variety;
  non-degeneracy of an irreducible component of \(\rho_{\mathcal{A}}^{-1}(X)\) follows from
  \cref{lem:betti_rank_zar_open_app} and the local homeomorphism property of the isogeny.
  \textit{Step (v):} Replace \(\mathcal{A}' \to S'\) by the principally polarized abelian scheme
  \(\mathcal{A}'_0 \to S'\) isogenous to it; non-degeneracy lifts because the isogeny is
  a local homeomorphism on the Betti torus.
  \textit{Step (vi):} Pass to a finite étale cover \(S'' \to S'\) on which \(\mathcal{A}'' = \mathcal{A}'_0 \times_{S'} S''\)
  acquires level-\(\ell\)-structure (via \cref{prop:betti_map_app}); this brings us into
  the setting of \cref{thm:ht_inequality}, which now applies to give the desired lower bound.
  Descent back along all the covers via the Height Machine yields the stated inequality on a
  Zariski open dense subset of the original \(X\).
\end{proof}
